% Part: incompleteness
% Chapter: representability-in-q
% Section: beta-function

\documentclass[../../../include/open-logic-section]{subfiles}

\begin{document}

\olfileid[pt]{inc}{req}{bet}
\olsection{O lema da função beta}


Para mostrar que podemos efetuar recursão primitiva quando dispomos de
adição, multiplicação e $\Char{=}$, precisamos desenvolver funções que
tratem sequências. (Se também tivéssemos exponenciação, nossa tarefa
seria mais fácil.) Quando dispúnhamos de recursão primitiva, podíamos
definir coisas como o ``$n$-ésimo primo'' e escolher uma codificação
bastante direta. Aqui, porém, não dispomos de recursão primitiva---na
verdade, queremos mostrar que podemos realizá-la usando minimização---,
e por isso precisamos ser mais engenhosos.

\begin{lem}
\ollabel{lem:beta}
Existe uma função $\beta(d,i)$ tal que, para toda sequência $a_0$,
\dots,~$a_n$, existe um número~$d$ tal que, para todo $i \le n$,
$\beta(d,i) = a_i$. Além disso, $\beta$ pode ser definida a partir das
funções básicas usando apenas composição e minimização regular.
\end{lem}

Pense em $d$ como um código da sequência $\tuple{a_0, \dots, a_n}$ e
em $\beta(d,i)$ como a função que devolve o $i$-ésimo elemento. (Note
que essa ``codificação'' \emph{não} usa a codificação por potências de
primos que já conhecemos!) O lema é bastante minimalista: não afirma
que podemos concatenar sequências ou acrescentar elementos, nem sequer
que podemos \emph{computar}~$d$ a partir de $a_0$, \dots,~$a_n$ usando
funções definíveis por composição e minimização regular. Afirma apenas
que existe uma função de ``decodificação'' tal que toda sequência é
``codificada''.

O uso da notação $\beta$ é de G\"odel. Repetindo: a parte difícil da
prova do lema é definir uma~$\beta$ adequada usando recursos
aparentemente restritos, O.s.\@, apenas composição e minimização---mas
podemos usar adição, multiplicação e~$\Char{=}$. Há várias maneiras de
provar esse lema, mas uma das mais elegantes continua sendo o método
original de G\"odel, que usava um fato da teoria dos números chamado
Teorema de Sunzi (tradicionalmente, o ``Teorema Chinês do Resto'').

\begin{defn}
Dois números naturais $a$ e $b$ são \emph{primos entre si} sse seu
máximo divisor comum é~$1$; em outras palavras, não têm nenhum outro
divisor em comum.
\end{defn}

\begin{defn}
Os números naturais $a$ e $b$ são \emph{congruentes módulo~$c$},
$a \equiv b \mod c$, sse $c \mid (a-b)$, O.s.\@, $a$ e $b$ deixam o
mesmo resto quando divididos por~$c$.
\end{defn}

Eis o Teorema de Sunzi:
\begin{thm}
Suponha que $x_0$, \dots,~$x_n$ sejam dois a dois primos entre si.
Sejam $y_0$, \dots,~$y_n$ números quaisquer. Então existe um número $z$
tal que
\begin{align*}
z & \equiv y_0 \mod x_0 \\
z & \equiv y_1 \mod x_1 \\
& \vdots  \\
z & \equiv y_n \mod x_n.
\end{align*}
\end{thm}

Usaremos o Teorema de Sunzi da seguinte maneira: se $x_0$,
\dots,~$x_n$ forem maiores, respectivamente, que $y_0$, \dots,~$y_n$,
podemos tomar $z$ como código da sequência $\tuple{y_0, \dots,y_n}$.
Para recuperar~$y_i$, basta dividir $z$ por~$x_i$ e tomar o resto. Para
usar essa codificação, precisaremos encontrar valores adequados para
$x_0$, \dots,~$x_n$.

Algumas observações nos ajudarão nesse ponto. Dados
$y_0$, \dots,~$y_n$, sejam
\begin{align*}
j &= \max(n, y_0 + 1, \dots, y_n + 1), \\
m &= \lcm(1,\dots,j),
\end{align*}
e sejam
\begin{align*}
x_0 & = 1 + m \\
x_1 & = 1 + 2 \cdot m \\
x_2 & = 1 + 3 \cdot m \\
& \vdots  \\
x_n & = 1 + (n+1) \cdot m
\end{align*}
Então duas coisas são verdadeiras:
\begin{enumerate}
\item\ollabel{rel-prime} $x_0,\dots,x_n$ são primos entre si.
\item\ollabel{less} Para cada $i$, $y_i < x_i$.
\end{enumerate}
Para ver que \olref{rel-prime} é verdadeira, observe que, se $p$ é um
número primo, $p \mid x_i$ e $p \mid x_k$, então
$p \mid 1 + (i+1) m$ e $p \mid 1 + (k+1) m$. Mas, nesse caso, $p$
divide a diferença entre eles,
\[
(1 + (i+1)m) - (1+ (k+1)m) = (i-k) m.
\]
Como $p$ divide $1 + (i+1)m$, ele não pode dividir também $m$ (caso
contrário, a primeira divisão deixaria resto~$1$). Logo, $p$ divide
$i-k$, pois $p$ divide $(i-k)m$. Mas $\left|i-k\right|$ é no máximo~$n$
e escolhemos $j \geq n$; portanto, isso implica que $p \mid m$, outra
contradição. Assim, não há número primo que divida simultaneamente
$x_i$ e $x_k$. A cláusula~\olref{less} é fácil: temos
$y_i < j \leq m < x_i$.

Provemos agora o lema da função $\beta$. Lembre-se de que podemos usar
$0$, sucessor, adição, multiplicação, $\Char{=}$, projeções e qualquer
função definida a partir deles usando composição e minimização aplicada
a funções regulares. Também podemos usar uma relação se sua função
característica for definível dessa forma. Como antes, podemos mostrar
que essas relações são fechadas sob combinações booleanas e
quantificação limitada; por exemplo:
\begin{align*}
\fn{not}(x) & \defis \Char{=}(x,0)\\
\bmin{x \leq z}{R(x,y)} & \defis \umin{x}{(R(x,y) \lor x = z)}\\
\bexists{x \leq z}{R(x,y)} & \defiff R(\bmin{x \leq z}{R(x,y)}, y)
\end{align*}
Podemos então mostrar que todos os itens a seguir também são definíveis
sem recursão primitiva:
\begin{enumerate}
\item A função de emparelhamento,
  $J(x,y) = \frac{1}{2}[(x+y)(x+y+1)] + x$;
% talvez explicar melhor o que está acontecendo aqui; é um pouco confuso.
\item as funções de projeção
\begin{align*}
K(z) & = \bmin{x \leq z}{\bexists{y \leq z}{z = J(x,y)}},\\
L(z) & = \bmin{y \leq z}{\bexists{x \leq z}{z = J(x,y)}};
\end{align*}
\item a relação de menor que $x < y$;
\item a relação de divisibilidade $x \mid y$;
% \item $x \tsub y$
% \item $\fn{Prime}(x)$
% \item Assuming $p$ is prime, the relation ``$x$ is a power of $p$'':
% \[
% \bforall{y \leq x}{(y \mid x \lif y = 1 \lor y = x)}.
% \]
\item a função $\fn{rem}(x,y)$, que devolve o resto da divisão de $y$
  por~$x$.
\end{enumerate}
Defina agora
\begin{align*}
\beta^*(d_0,d_1,i) & = \fn{rem}(1+(i+1) d_1,d_0) \text{ e}\\
\beta(d,i) & = \beta^*(K(d),L(d),i).
\end{align*}
Esta é a função que queremos. Dados $a_0,\dots,a_n$ como acima, seja
\[
j = \max(n,a_0+1,\dots,a_n+1),
\]
e seja $d_1 = \lcm(1,\dots,j)$. Por \olref{rel-prime}, sabemos que
$1+d_1$, $1+2 d_1$, \dots, $1+(n+1) d_1$ são primos entre si e, por
\olref{less}, que todos são maiores que $a_0,\dots,a_n$. Pelo Teorema
de Sunzi, existe um valor~$d_0$ tal que, para cada~$i$,
\[
d_0 \equiv a_i \mod (1+(i+1)d_1)
\]
e, portanto (como $d_1$ é maior que~$a_i$),
\[
a_i = \fn{rem}(1+(i+1)d_1,d_0).
\]
Seja $d = J(d_0,d_1)$. Então, para cada $i \le n$, temos
\begin{align*}
\beta(d,i) & =  \beta^*(d_0,d_1,i) \\
& =  \fn{rem}(1+(i+1) d_1,d_0) \\
& =  a_i
\end{align*}
que é o que precisávamos. Isso conclui a prova do lema da função
$\beta$.

\begin{prob}
  Mostre que as relações $x < y$, $x \mid y$ e a
  função~$\fn{rem}(x,y)$ podem ser definidas sem recursão primitiva.
  Você pode usar $0$, sucessor, adição, multiplicação, $\Char{=}$,
  projeções, minimização limitada e quantificação limitada.
\end{prob}

\end{document}
